\chapter{Pruebas y evaluación}
\label{cap:evaluacion-resultados}

\section{Pruebas disponibles}
En la última ejecución completa registrada (13/11/2026) terminaron
correctamente 119 pruebas de backend (\texttt{mvnw verify}: gateway 1,
identity 12, election 46, voter-access 38 y ballot-box 22) y 59 de frontend
(\texttt{npm run test:ci}); \texttt{npm run build} también terminó
correctamente. Las pruebas unitarias y de arquitectura cubren reglas
seleccionadas de permisos, ciclo de vida, importación del censo, acceso del
votante, consumo del token, cadena de auditoría, escrutinio y separación
arquitectónica. En frontend se cubren la guarda de rol, el flujo de voto y
recibo, el retorno SSO con PKCE, componentes del panel de la consola, el
asistente de creación de procesos, la traducción automática de sus textos y el
aviso de caducidad de la sesión.

Dieciséis de las pruebas de backend, repartidas en ocho clases, se ejecutan
contra PostgreSQL 18.6 con Testcontainers. Aplican las migraciones Flyway de
identity, election, voter-access y ballot-box y verifican el trigger
append-only de acciones administrativas, la participación con fecha sin hora
y las restricciones de esquema y eventos de la urna. Además comprueban que una
organización no puede leer ni modificar los datos de otra en election,
voter-access y ballot-box; que el despachador del \emph{outbox} reintenta
tras un fallo y entrega una sola vez; que veinte intentos simultáneos con el
mismo token producen exactamente un voto; que ninguna tabla de voter-access
puede almacenar la opción elegida ni un recibo; y que la emisión no escribe
en los registros el token, el recibo ni la opción. No cubren la autorización
HTTP de identity ni todos los repositorios y entidades.

Los barridos de accesibilidad con Playwright y axe de esa misma ejecución (200
comprobaciones: 108 de la consola, 88 del flujo del votante y 4 de la
comprobación del recibo; 199 sin incumplimientos críticos ni serios y uno con
un incumplimiento serio en la consola, en la edición de la papeleta a 360
píxeles: tras avanzar de paso, la barra de pasos queda bajo la cabecera, tapada
por el selector de idioma; corregido el mismo día haciendo que el asistente
vuelva a la barra de pasos al cambiar de paso, y repetidas las 200
comprobaciones sin incumplimientos críticos ni serios) se valoran en la
tabla~\ref{tab:estado-hipotesis}; no son pruebas funcionales de extremo a
extremo.

La cobertura de líneas JaCoCo de esa ejecución es: gateway 33,3\% (1/3),
identity 14,4\% (113/785), election 70,3\% (1267/1801), voter-access 50,2\%
(789/1572) y ballot-box 41,1\% (672/1636). En el frontend, medida con Vitest
sobre todo el código de \texttt{src/app} (14/11/2026), es del 22,7\% de las líneas
(486/2139). La cobertura no demuestra por sí
misma la corrección del sistema. El informe de evaluación detalla método y
limitaciones en \texttt{docs/evaluation/informe.md}.

\section{Hipótesis y resultados}
\small
\begin{longtable}{|p{0.11\textwidth}|p{0.36\textwidth}|p{0.15\textwidth}|p{0.26\textwidth}|}
\caption{Estado de las hipótesis de evaluación.}\label{tab:estado-hipotesis}\\
\hline
\textbf{Hipótesis} & \textbf{Criterio del protocolo} & \textbf{Estado del criterio} & \textbf{Resultado medido} \\
\hline
\endfirsthead
\hline\textbf{Hipótesis} & \textbf{Criterio del protocolo} & \textbf{Estado del criterio} & \textbf{Resultado medido} \\
\hline\endhead
H1a \mbox{(RNF-1)} & Cero incumplimientos críticos ni serios en las pantallas evaluadas, más revisión manual con teclado y lector de pantalla sin bloqueos en el flujo de acceso, voto y recibo. La auditoría automática no certifica conformidad WCAG. & Validado por el tutor & parte automática favorable; parte manual no realizada. Barrido axe del flujo del votante (05/11/2026; 22 pantallas o estados, 1280 y 360 píxeles, español e inglés; 88 auditorías): el primer barrido dio 8 incumplimientos críticos de una misma incidencia en la papeleta (\texttt{aria-required-children}); corregida, el segundo barrido da 0 críticos y 0 serios. También se corrigió una incidencia que axe no detecta (el error de formato del código de recibo no se mostraba); el tercer barrido, con las dos correcciones, da 0 críticos y 0 serios. Repetido el 08/11/2026 sobre la versión vigente (\texttt{06f39be}), con el mismo resultado. La revisión manual con teclado y lector de pantalla que exige el criterio no se realizó antes del cierre: no es un incumplimiento, sino una parte del protocolo sin ejecutar. Evidencia parcial; sin decisión. \\
\hline
H1b \mbox{(RNF-3,} \mbox{RNF-6)} & Los 5 participantes no técnicos completan acceso, emisión y obtención del recibo en $\leq$ 3 minutos, y al menos 4 de 5 interpretan correctamente el resultado de la verificación. H0b: no se alcanza alguno de esos criterios. & Validado por el tutor & ACEPTADA: sesiones del 09/11/2026 con consentimiento y escenario sintético. Cinco sesiones válidas (P02--P05 y la reserva R1; P01 se invalidó por una incidencia de entorno, porque su entrada del censo ya tenía un voto). 5 de 5 completan acceso, emisión y recibo sin ayuda en $\leq$ 3 minutos (T0--T1: mediana 51~s, mínimo 36~s, máximo 79~s; 0 errores) y 5 de 5 interpretan correctamente el recibo. Una sesión adicional con la segunda reserva (R2) queda fuera del cómputo y tuvo el mismo resultado. Evidencia descriptiva limitada al prototipo, la configuración y la población evaluados. \\
\hline
H1c & Con 500 usuarios simultáneos, p95 de emisión inferior a dos segundos en tres ejecuciones. & Propuesto & RECHAZADA (H0c): tres ejecuciones del 11/11/2026 (procesos sintéticos 10, 11 y 12; 500 usuarios virtuales que empiezan a la vez, un voto cada uno). p95 de la emisión: 4852, 2402 y 4849~ms, por encima de 2000~ms en las tres. Votos aceptados: 222, 240 y 281 de 500; el resto fueron conexiones rechazadas por el \emph{gateway} (273, 260 y 219) y, en la primera, 5 errores 500. En las tres, la cadena quedó íntegra y los votos de la urna coinciden con las respuestas 201 y con la participación. Un único equipo Windows en local; sin extrapolación a producción. \\
\hline
RNF-4 (sin hipótesis) & Al menos 4 de 5 administradores de organización sin formación previa crean, configuran y programan un proceso de opción única en $\leq$ 10 minutos y sin ayuda. & Validado por el tutor; 10 minutos como valor propuesto & ACEPTADO: sesiones online del 13/11/2026 (videollamada grabada, consentimiento escrito, escenario sintético), cinco válidas y las cinco sin formación previa en VeriVoto. 5 de 5 crean, configuran y programan el proceso sin ayuda en $\leq$ 10 minutos (mediana 3 min, mínimo 2 y máximo 4; 0 errores; ninguno consultó la guía «Primeros pasos»). Tiempos medidos por el moderador sobre la pantalla compartida y anotados redondeados al minuto; con un máximo de 4 minutos frente a 10, el redondeo no cambia la decisión. Sustituye a la sesión del 03/11/2026, no concluyente (solo 2 de 5 sin formación previa y tiempos estimados). \\
\hline
\end{longtable}
\normalsize

Con $n = 5$ no hay inferencia estadística: los resultados de H1b y de RNF-4
son evidencia descriptiva limitada al prototipo, la configuración y la
población evaluados. La columna «Resultado medido» solo se completa con
ejecuciones reales.

En H1c, el p95 de k6 incluye las conexiones rechazadas, que terminan casi al
instante; el de los votos aceptados es todavía mayor (5055, 2477 y 5228~ms).
Las conexiones se rechazaron en el \emph{gateway} al llegar 500 a la vez y, en
la primera ejecución, también entre el \emph{gateway} y ballot-box. El registro
del voto bloquea la fila del proceso para encadenar cada evento de la urna, de
modo que los votos de un mismo proceso se guardan de uno en uno; es una
explicación compatible con los tiempos observados, pero no se ha medido dónde
se consume el tiempo. El criterio no se modifica a la vista del resultado.

\section{Integración y requisitos de calidad}
La integración PostgreSQL de otros repositorios, la autorización HTTP de
identity, las pruebas funcionales de extremo a extremo del stack completo y la
accesibilidad manual continúan pendientes; la sesión con votantes (H1b) se realizó el 09/11/2026 y
las tres ejecuciones de carga (H1c), el 11/11/2026. En consecuencia, la
evidencia disponible no permite declarar aceptados RNF-1 ni
RNF-11--RNF-15 en su conjunto, y RNF-17 no se cumple en la configuración
local evaluada; para RNF-3 y RNF-6 se dispone de la evidencia
descriptiva de H1b, con cinco participantes, y RNF-4 se acepta con la de las
sesiones con administradores del 13/11/2026.
La correspondencia con características de ISO/IEC 25010:2023 está definida
en el capítulo de requisitos no funcionales; la cobertura experimental por
característica debe completarse cuando se ejecuten esos protocolos.

\section{Validez}
La cobertura es desigual (del 14,4\% de identity al 70,3\% de election) y
limita la confianza que puede atribuirse a la suite. Las pruebas de
integración comprueban casos concretos de aislamiento entre organizaciones y
de concurrencia sobre una única instancia; no sustituyen la verificación de
todos los repositorios y rutas ni la concurrencia real de un despliegue. El rechazo de H1c se limita a la configuración
evaluada (un único equipo Windows, servicios fuera de contenedores y 500
conexiones simultáneas sin rampa) y no permite estimar el rendimiento de un
despliegue dimensionado, y el resultado de H1b, obtenido con cinco participantes
autoseleccionados en un entorno local, no se extiende a otros votantes,
dispositivos ni configuraciones. Los resultados aquí consignados se
limitan a las ejecuciones disponibles y no se extrapolan a producción.
La parte manual de H1a (teclado y lector de pantalla) no se ejecutó: el
barrido automático solo cubre las reglas que axe puede decidir y dejó
avisos sin resolver (contraste sobre fondos que no puede calcular y un
atributo ARIA del recibo), y no comprueba cómo se anuncian los errores ni el
orden del foco. Por eso H1a no se declara aceptada ni rechazada y queda como
evidencia parcial y sin decisión.

El resultado de RNF-4 tiene sus propias limitaciones: las cinco personas son
estudiantes autoseleccionados del entorno del autor; la tarea se hizo a
distancia, con un código de acceso fijo para estas cuentas (el inicio de
sesión queda fuera del tiempo medido); los tiempos se anotaron redondeados al
minuto; tres participantes escribieron textos propios en lugar de los de la
ficha, y dos usaron la versión de la consola anterior a los ajustes aplicados
durante el bloque de sesiones a partir de las observaciones de los primeros
participantes. Nada de ello afecta a la tarea medida (crear, configurar y
programar un proceso de opción única), pero limita la generalización.
